\section{Normalization and the arithmetic comparison}
\label{sec:normalization}

\subsection{A precise canonical syntax}

For the metatheory, first distinguish the raw constructors $0,S,R$.
Orient the two recursor equations as
\[
 R(0,a;i,u.h)\longrightarrow Sa,\qquad
 R(Sn,a;i,u.h)\longrightarrow h[n/i,R(n,a;i,u.h)/u].
\]
Reductions may occur in every numerical subterm, including an update
body under its two binders. The mathematical normal form is taken in
this raw syntax. Its finite encoding then prints $Sa$ as $[0,a,0]$ and
an unreduced $R(n,a;i,u.h)$ as $[n,a,h]$. Thus the implementation's
zero-counter clause is a canonical constructor operation, not a rewrite
of $R(0,a;0)$ to itself to be repeated forever.

\begin{lemma}[Canonical computation]
Every well-sorted finite term has a unique computation normal form.
Reduction introduces no new free variable. For every capture-avoiding
substitution $\sigma$,
\[
 N(t[\sigma])=N(N(t)[N\circ\sigma]).
\]
The corresponding statements hold for formulas and one-place templates.
\end{lemma}
\begin{proof}
One standard strong-normalization proof embeds a raw recursor into the
ground-type natural-number recursor: its starting value is $Sa$, and its
step is the finite function $(i,u)\mapsto h(i,u)$. Variables in $h$ are
bound numerical arguments; there is no application of a variable as a
function. The usual reducibility proof specializes directly to this
fragment. Its recursor case inducts on a computable counter and its
reducts: at zero the result is the computable starting value, and at a
successor the recursive call is computable at the predecessor and the
computable update preserves computability. At a neutral counter the
recursor is neutral and its reductions are those of computable proper
subterms. Structural induction on a finite term, with computable terms
substituted for its numerical variables, establishes computability and
hence strong normalization. Applying the same argument under fresh
variables handles reductions inside update bodies.

The two root patterns are disjoint and left-linear. Nested reductions
commute by the capture-avoiding substitution lemma; duplicating an update
or an argument duplicates its reducts on the other branch of the diagram.
The parallel-reduction argument therefore gives confluence. Strong
normalization and confluence give a unique normal form. Applying a
substitution to each computation equation yields another valid
computation; normalizing the two sides of this conversion proves the
displayed coherence identity. All right sides use only variables from
their left sides, proving the free-variable assertion.
\end{proof}

This is an external metatheorem about finite expressions. There is no
claim that normalization of arbitrary supplied descriptions admits a
uniform primitive-recursive bound. In particular, a full normalizer is
not free infrastructure in the complexity comparison: its whole source
is included, and the bounded verifier charges every recursive traversal.

\subsection{Computation is also provable in the ordinary theory}

The ordinary realization $T_R$ mentioned above uses raw finite recursor
descriptions. An instance nested in a counter or initializer is transported
by equality. A reduction inside an update is justified by the following
body-extensionality lemma, rather than by silently treating a bound
function as a first-order numerical argument.

\begin{lemma}[Extensionality for a fixed update]
If $h(i,u)=k(i,u)$ is derivable uniformly in its two displayed variables,
then
\[
 R(n,a;i,u.h)=R(n,a;i,u.k)
\]
is derivable for every $n$, with all passive parameters fixed.
\end{lemma}
\begin{proof}
Use literal-equation induction on the displayed equality. At zero both
sides are $Sa$. At a successor, computation exposes the two updates.
The induction equation first replaces the old state in one update;
the supplied pointwise update identity then changes $h$ into $k$.
All inductions used here have a literal equation as their formula.
\end{proof}

Consequently each finite reduction to $N(t)$ yields an internal proof
of $t=N(t)$. An external induction on a formula gives the corresponding
equivalence between its raw and normalized readings. There is no use
of a uniform internal truth or evaluation predicate.

\subsection{The full comparison is over finite descriptions}

The comparison theory for the proof translation is the diagonal native
calculus, conservatively supplied with fixed $R$ descriptions and their
computation equations. Each occurrence of $R(n,a;i,u.h)$ first compiles
the proper subdescriptions of $a,h$ and then introduces the primitive
recursion with that fixed update. The resulting description is finite.
This construction is staged by syntax, not by a runtime code variable.

The diagonal proof works uniformly in the enlarged numerical vocabulary:
its selector identity and its equation-induction simulation treat the two
sides of an equation as arbitrary term templates. The new computations
translate to literal reflexivity after normalization. Conversely, iPRA
interprets every such fixed description by primitive recursion. This
supplies the whole target language, including equations with $R$ in their
updates, to the source-to-target argument.

For clarity, the forward compilation of an ordinary primitive-recursive
definition is explicit:
\[
\begin{aligned}
0^*&=0, & S^*(x)&=R(0,x;i,u.0), & \pi_j^*(\vec x)&=x_j,\\
(f\circ\vec g)^*(\vec x)&=f^*(g_1^*(\vec x),\ldots,g_m^*(\vec x)),\\
(\operatorname{rec}(g,h))^*(n,\vec x)
  &=P\bigl(R(n,g^*(\vec x);i,v.
        S(h^*(i,P(v),\vec x)))\bigr).
\end{aligned}
\]
The reverse compilation sends $R(n,a;i,u.h)$ to the fixed ordinary
primitive-recursive function with base $S(a)$ and step $h$. These are
translations of descriptions; no quantification over functions is added.

\begin{theorem}[Bi-interpretability and consistency]
The ordinary realization of the three-instruction calculus is
bi-interpretable with $\ipra$. Both domains are the full numerical sort,
equality is genuine equality, and both numerical comparison maps have
graph $x=y$. Each full formula is provably equivalent to its double
translation. In particular the new calculus proves $0=S0$ if and only
if $\ipra$ proves $\bot$.
\end{theorem}
\begin{proof}
The rule-by-rule argument establishes the presentation theorem for
$T_R$. The interpretation of its calculations in $\ipra$ is the fixed
recursion just described. Coverage and separation hold in $\ipra$, and
literal-equation induction is an instance of quantifier-free induction.
For the reverse direction the displayed compilation validates the
primitive-recursive equations; applying $P$ proves successor injectivity.
If $Sx=0$, the selector equations force $0=S0$, so successor separation
also holds. The scalar iterator $J$ supplies the already established
scalar equality-reflection construction. Thus every quantifier-free
formula is equivalent to the vanishing of a fixed numerical term, and
its induction instance follows from equation induction. The attributed
full equality-reflection derivation is reproduced in the reference folder
\cite{foundations}.

For a fixed recursive definition, equation induction proves uniqueness
from its base and step equations. External induction on each finite
description therefore proves that its double translation has the same
value as the original, internally for every numerical argument and every
passive parameter. The domain interpretations are unrestricted and the
isomorphism is the identity. The native logical encodings below and the
ordinary reading of $D$ supply both formula round trips. Falsehood is
carried to the fixed equation $0=S0$, and its ordinary reading implies
$\bot$ by the explicit separation axiom. Hence contradiction proofs
translate both ways.
\end{proof}

The argument asserts effective finite proof translations and internal
round-trip identities. It does not assert a uniform primitive-recursive
normalization algorithm, a proof of consistency inside iPRA, or a new
formal verification of the metatheory.

\subsection{Recovering the ordinary logical vocabulary}

Here is an explicit account of the logical language being recovered.
Write $q_yA=D_{d,y}A$ with $d$ unused. Formula padding is
$L(A)=q_yG_yA$, where $G$ gates the final equation by $C$. Padding is
equivalent to $A$; it permits two formulas to be brought to the same
number of paired binders. After freshening and aligning those binders,
replace their two final equation endpoints componentwise by the selector
$C(v,-,-)$. The resulting formula $\operatorname{sel}(v,A,B)$ is
equivalent to $A$ at zero and $B$ at a successor. All temporary bound
variables are distinct from $v$ and the free variables of $A,B$.

Let $\operatorname{par}$ be the fixed primitive-recursive parity term.
One standard native encoding is
\[
\begin{aligned}
\bot^*&=(0=S0),\\
(A\to B)^*&=q_y\operatorname{sel}(y,A^*,B^*),\\
(A\land B)^*&=q_yG_y\operatorname{sel}(\operatorname{par}(y),A^*,B^*),\\
(\exists x\,A)^*&=D_{x,y}G_yA^*,\\
(A\lor B)^*&=\exists v\,\operatorname{sel}(v,A^*,B^*),\\
(\forall x\,A)^*&=A^*(0)\land q_x A^*(x).
\end{aligned}
\]
The last three displays use the macros defined in earlier lines.
Bound variables in each line are fresh. Selecting $1$ and $2$ in the
conjunction gives its two components. For implication, its zero row is
the antecedent and each successor row is the consequent. Existential
introduction and elimination are exactly the gated paired rules. Numerical
coverage proves the disjunction case rule. Transport along $0=S0$ in
a gate changes its provable zero instance into any chosen formula,
giving explosion. These are intuitionistic derivations; no excluded
middle for formulas with quantifiers is used.

The diagonal native equivalence gives these rules under arbitrary retained
assumptions, not just for closed arithmetic identities. By structural
induction the ordinary reading of each macro is equivalent to its intended
connective. Expanding the ordinary reading of a native $D$ back through
the macros gives the original $D$, again under arbitrary assumptions.
Together with the numerical uniqueness argument this proves both formula
round trips, rather than only mutual interpretability of theorem sets.

